\documentclass[12pt,a4paper]{article}

% Drake_preamble.tex - LaTeX Preamble Template
% 作者: 謝慕揚, MD, PhD, FESC

% Including essential packages for Chinese typesetting
\usepackage{xeCJK}
\usepackage{fontspec}
% Configuring Chinese font with Noto Serif CJK TC, fallback to SimSun
\IfFontExistsTF{Noto Serif CJK TC}{
  \setCJKmainfont{Noto Serif CJK TC}
  \setCJKsansfont{Noto Serif CJK TC}
  \setCJKmonofont{Noto Serif CJK TC}
}{\setCJKmainfont{SimSun}
  \setCJKsansfont{SimSun}
  \setCJKmonofont{SimSun}
}

% Including other necessary packages
\usepackage{geometry}
\usepackage{titlesec}
\usepackage{enumitem}
\usepackage{xcolor}
\usepackage{colortbl}
\usepackage{tcolorbox}
\usepackage{booktabs}
\usepackage{longtable}
\usepackage{array}
\usepackage{graphicx}
\usepackage{tikz}
\usepackage{fancyhdr}
\usepackage{multicol}
\usepackage{datetime2}
\usepackage{hyperref}
\usepackage{mdframed}
\usepackage{amsmath}
\usepackage{amssymb}
\usepackage{multirow}

% Defining colors
\definecolor{emergency_red}{RGB}{186,24,27}
\definecolor{emergency_blue}{RGB}{0,114,188}
\definecolor{emergency_gray}{RGB}{120,120,120}
\definecolor{highlight_yellow}{RGB}{255,250,205}

\hypersetup{
    colorlinks=true,
    linkcolor=emergency_blue,
    citecolor=emergency_blue,
    urlcolor=emergency_blue,
    pdfborder={0 0 0}
}

% Defining page geometry
\geometry{
  top=2.5cm,
  bottom=2.5cm,
  left=2cm,
  right=2cm,
  headheight=25pt
}

% Adjusting topmargin to compensate for increased headheight
\addtolength{\topmargin}{-10pt}

% Setting title formats
\titleformat{\section}[block]{\Large\bfseries\color{emergency_red}}{\thesection}{10pt}{}
\titleformat{\subsection}[block]{\large\bfseries\color{emergency_blue}}{\thesubsection}{8pt}{}
\titleformat{\subsubsection}[block]{\normalsize\bfseries\color{emergency_gray}}{\thesubsubsection}{6pt}{}

% Defining custom environment for important notes
\newmdenv[
    linecolor=emergency_red,
    backgroundcolor=highlight_yellow,
    linewidth=2pt,
    topline=true,
    bottomline=true,
    leftline=true,
    rightline=true,
    innertopmargin=10pt,
    innerbottommargin=10pt,
    innerleftmargin=10pt,
    innerrightmargin=10pt
]{importantbox}

% Defining note command with tcolorbox
\newcommand{\note}[1]{
    \par
    \noindent
    \begin{tcolorbox}[
        colback=emergency_blue!5!white,
        colframe=emergency_blue,
        title=\textbf{重點提醒},
        fonttitle=\bfseries,
        boxrule=0.5pt,
        arc=3pt,
        left=6pt,
        right=6pt,
        top=6pt,
        bottom=6pt
    ]
    #1
    \end{tcolorbox}
}

% Key findings box
\newcommand{\keyfinding}[1]{
    \par
    \noindent
    \begin{tcolorbox}[
        colback=yellow!10!white,
        colframe=orange!75!black,
        title=\textbf{核心發現摘要},
        fonttitle=\bfseries,
        boxrule=0.5pt,
        arc=3pt
    ]
    #1
    \end{tcolorbox}
}

% Defining custom date format
\newcommand{\mydate}{\the\year-\the\month-\the\day}

\begin{document}

%% Title Page
\begin{center}
{\LARGE\bfseries\textcolor{emergency_red}{EuroIntervention Expert Review}}\\[0.5cm]
{\Large\textbf{Endomyocardial Biopsy: A Practical Guide for Interventional Cardiologists}}\\[0.3cm]
{\large 心內膜心肌切片：介入心臟科醫師實務指南}\\[0.5cm]
{\normalsize EuroIntervention. 2026;21:e19-e31}\\[0.3cm]
{\normalsize 整理：謝慕揚 MD, PhD, FESC \quad 日期：\mydate}
\end{center}

\vspace{0.5cm}

\keyfinding{
\begin{itemize}[leftmargin=*]
    \item Endomyocardial biopsy (EMB) 是診斷 myocarditis、cardiomyopathies、storage/infiltrative diseases 的 gold standard
    \item 主要併發症發生率約 1-5\%，包括 cardiac perforation、thromboembolism、valvular trauma、sustained VT
    \item RV EMB 適用於 diffuse diseases（如 CA、Anderson-Fabry disease）；LV EMB 適用於 LV-predominant diseases（如 myocarditis、DCM）
    \item EMB 在 unexplained acute HF with hemodynamic compromise 中扮演關鍵角色，可達 78\% 診斷率
    \item Heart Team 討論及適當的 patient selection 是成功執行 EMB 的關鍵
\end{itemize}
}

\section{EMB 的臨床適應症}

\subsection{ESC Guidelines 建議}

根據 ESC Heart Failure Guidelines，EMB 適用於 (Class IIa)：
\begin{itemize}
    \item 「Rapidly progressive HF despite standard therapy, when there is a clinical pretest probability of a specific diagnosis requiring precise treatments, which can be confirmed only in myocardial samples」
\end{itemize}

\subsection{主要臨床適應症}

\begin{longtable}{p{5cm}p{9cm}}
\toprule
\textbf{臨床情境} & \textbf{說明} \\
\midrule
\endfirsthead
\toprule
\textbf{臨床情境} & \textbf{說明} \\
\midrule
\endhead
Suspected fulminant/acute myocarditis & 合併 acute HF、rhythm disorders 或 hemodynamically stable 但高度懷疑 \\
\midrule
Unexplained acute HF & 合併 hemodynamic compromise 或 ventricular arrhythmias/conduction disorders \\
\midrule
DCM with new-onset HF & LV dysfunction 對 standard medical therapy 無反應 \\
\midrule
Unexplained hypertrophic/restrictive phenotype & 需鑑別 sarcomeric HCM 與 phenocopies（如 CA） \\
\midrule
Suspected cardiac sarcoidosis & Focal disease，建議 biventricular EMB \\
\midrule
Storage/infiltrative diseases & Amyloidosis、Anderson-Fabry disease 等 \\
\midrule
ICI-mediated cardiotoxicity & Immune checkpoint inhibitor 相關 myocarditis \\
\midrule
Heart transplant rejection & 排斥反應監測 \\
\bottomrule
\end{longtable}

\subsection{EMB 不適用情境}

\begin{itemize}
    \item \textbf{Low-risk acute myocarditis}：僅有 chest pain、normal ventricular function、無 ventricular arrhythmias、ECG 異常快速緩解
    \item \textbf{Friable intracardiac masses}：高 embolic risk，如 left-sided tumours 或 typical cardiac myxomas
    \item \textbf{Mobile and small protruding structures}：技術上困難且不安全
\end{itemize}

\note{
EMB 具有特殊價值的情境：(i) acute cardiac settings refractory to standard therapies；(ii) 無法接受 CMR 的患者；(iii) heart transplant rejection surveillance；(iv) chronic stable patients with inconclusive non-invasive findings 且懷疑 inflammatory diseases。
}

\section{RV vs LV EMB 的選擇}

\subsection{選擇原則}

選擇 EMB approach 需考慮三大因素：
\begin{enumerate}
    \item \textbf{Clinical query}：臨床問題為何
    \item \textbf{Most diseased areas}：哪個心室受影響較嚴重
    \item \textbf{Nature of cardiac process}：Focal vs diffuse disease
\end{enumerate}

\subsection{RV vs LV EMB 比較}

\begin{longtable}{p{4cm}p{5cm}p{5cm}}
\toprule
\textbf{項目} & \textbf{RV EMB} & \textbf{LV EMB} \\
\midrule
\endfirsthead
\toprule
\textbf{項目} & \textbf{RV EMB} & \textbf{LV EMB} \\
\midrule
\endhead
首選適應症 & Biventricular involvement & LV diseases with preserved RV \\
\midrule
技術難度 & 較低 & 較高 \\
\midrule
Cardiac perforation & 風險較高但較少致命 & 風險較低但更致命 \\
\midrule
Access site complications & Venous complications（如 hematoma） & Arterial complications（如 pseudoaneurysm） \\
\midrule
Embolization risk & 較低（pulmonary embolism） & 略高（stroke 或 systemic embolism） \\
\midrule
Conduction injury & 較高 AV block 風險 & 較低 AV block 風險 \\
\midrule
Valve damage & Tricuspid valve & Mitral valve \\
\bottomrule
\end{longtable}

\subsection{診斷效能}

根據 Chimenti et al. 研究（$>$4,000 patients）：
\begin{itemize}
    \item 當 LV 單獨或主要受影響時：LV EMB diagnostic yield 97.8\%，RV EMB 僅 53\%
    \item 當 RV 同時受影響時：RV EMB diagnostic yield 增至 96.5\%，LV EMB 98.1\%
\end{itemize}

\subsection{臨床建議}

\begin{longtable}{p{7cm}p{7cm}}
\toprule
\textbf{Prefer RV EMB} & \textbf{Prefer LV EMB} \\
\midrule
\endfirsthead
\toprule
\textbf{Prefer RV EMB} & \textbf{Prefer LV EMB} \\
\midrule
\endhead
Cardiac amyloidosis & Suspected myocarditis with primary LV involvement \\
\midrule
Anderson-Fabry disease & Unexplained DCM with isolated LV involvement \\
\midrule
Arrhythmogenic cardiomyopathy & Cardiac sarcoidosis（建議 biventricular） \\
\midrule
Unexplained restrictive cardiomyopathy & \\
\midrule
AM with biventricular abnormality & \\
\bottomrule
\end{longtable}

\section{技術與操作步驟}

\subsection{術前準備}

\begin{itemize}
    \item 確認無 ventricular thrombus
    \item Platelet count $>$50,000/mm$^3$
    \item 停用 oral anticoagulants
    \item Ultrasound-assisted puncture 減少 access site complications
    \item 持續監測 heart rhythm 及 invasive blood pressure
    \item 貼好 defibrillator patches
\end{itemize}

\subsection{基本操作步驟}

\begin{enumerate}
    \item 建立 vascular access（RV: jugular/femoral/brachial vein；LV: femoral/radial artery）
    \item LV EMB 需給予 IV heparin 達 ACT $>$200 seconds
    \item 經由 pigtail catheter over 0.035" guidewire 將 long sheath 送入心室
    \item 執行 ventriculography 確認位置
    \item 移除 pigtail，aspirate and flush sheath，確認 ventricular pressure
    \item 以 fluoroscopy 確認 sheath tip 位置（RAO + LAO projections）
    \item 將 prebent bioptome 送入 long sheath
    \item 在 sheath 遠端打開 forceps，保持 open 直到接觸心肌
    \item 感受輕微阻力時關閉 jaws（可能出現 ventricular ectopy 或 NSVT）
    \item 關閉並撤回為一個連續動作
    \item 撤回時保持 jaws closed 以避免 specimen loss 或 embolization
    \item 重複取樣至少 5 個不同位置
    \item 術後 echocardiography 確認無 pericardial effusion 或 valvular damage
\end{enumerate}

\note{
LV EMB 進行時，若可能應啟動 cardiac surgery standby。每位患者術後需在 cath lab 監測至少 15 分鐘並執行 TTE。
}

\subsection{Fluoroscopy Projections}

\begin{itemize}
    \item \textbf{RV EMB}：
    \begin{itemize}
        \item RAO 30° 用於確認 RV outflow tract 和 RV free wall 位置
        \item LAO 40° 最常用於引導 septum sampling（提供 interventricular septum profile view）
        \item 避免在 outflow tract 太近取樣以免損傷 right bundle branch
    \end{itemize}
    \item \textbf{LV EMB}：
    \begin{itemize}
        \item RAO + LAO projections 確認 mid-LV cavity 位置
        \item 避開 apex 並遠離 valvular apparatus
    \end{itemize}
\end{itemize}

\section{併發症}

\subsection{主要併發症（1-5\%）}

\begin{itemize}
    \item \textbf{Cardiac perforation}：最嚴重，RV perforation 較常見但 LV perforation 更致命
    \item \textbf{Thromboembolism}：LV EMB 風險較高（stroke/systemic embolism）
    \item \textbf{Valvular trauma}：Tricuspid（RV EMB）或 Mitral（LV EMB）
    \item \textbf{Sustained ventricular arrhythmias}：通常 transient
    \item \textbf{Vascular complications}：Access site hematoma、pseudoaneurysm
    \item \textbf{Conduction system injury}：RV EMB 較高風險
\end{itemize}

\subsection{特殊考量}

\begin{itemize}
    \item 患者有 LBBB 時，RV EMB 需準備 temporary pacemaker（complete AV block 風險）
    \item VA-ECMO 支持的 fulminant myocarditis 患者，EMB 可達 78\% diagnostic yield，但併發症風險顯著增加
\end{itemize}

\section{與 Multimodality Imaging 整合}

\subsection{Preprocedural Imaging}

\begin{itemize}
    \item \textbf{CMR + 18F-FDG PET}：改善 cardiac sarcoidosis 的 EMB 診斷準確度
    \item CMR 與 EMB 在 myocarditis 的 concordance 有限，支持兩者作為 complementary tools
\end{itemize}

\subsection{Procedural Imaging}

\begin{itemize}
    \item \textbf{Electroanatomical mapping (EAM)-guided EMB}：
    \begin{itemize}
        \item Low-voltage areas 與 histological abnormalities 相關
        \item 特別適用於 segmental/patchy involvement（如 cardiac sarcoidosis）
        \item 較耗時且成本較高
    \end{itemize}
    \item \textbf{Intracardiac echocardiography}：確認 bioptome positioning
\end{itemize}

\section{Tissue Processing 與分析技術}

\subsection{Myocarditis 診斷}

\begin{itemize}
    \item \textbf{Dallas criteria}（1980s）：sensitivity 低
    \item \textbf{Immunohistochemistry}：
    \begin{itemize}
        \item HLA-DR：immune activation
        \item CD3, CD4, CD8：T cells
        \item CD68：macrophages
    \end{itemize}
    \item \textbf{Viral PCR}：決定是否適合 immunosuppression
\end{itemize}

\subsection{Amyloidosis 診斷}

\begin{itemize}
    \item \textbf{Congo Red staining}：Apple-green birefringence under polarized light
    \item \textbf{Electron microscopy}：Non-branching 10 nm fibrils
    \item \textbf{Immunohistochemistry}：Detect TTR、light chains、apolipoproteins
    \item \textbf{Mass spectrometry}：Preferred technique for amyloid typing
\end{itemize}

\section{EMB 在 Cardiogenic Shock 的角色}

\begin{itemize}
    \item 在 non-ischemic severe acute HF/cardiogenic shock 且懷疑 fulminant myocarditis 時，EMB 可識別 underlying cause 並指導 targeted treatment
    \item 早期 EMB 與較佳的 1-year outcomes 相關
    \item 需要 mechanical support 但未做 EMB 者預後較差（可能因延遲診斷與治療）
    \item 應轉介至有經驗的專門中心執行
\end{itemize}

\note{
Immunosuppression 適應症：Giant cell myocarditis、eosinophilic myocarditis、granulomatous myocarditis（cardiac sarcoidosis）需及時啟動 immunosuppressive therapy 以改善 survival。對於 virus-negative lymphocytic myocarditis，個人化的 immunosuppression 是有效且安全的。
}

\section{Case Examples 摘要}

\subsection{Case 1：Cardiomyopathy with Cardiogenic Shock}

\begin{itemize}
    \item 32 歲男性，upper respiratory infection 後 1 個月出現 palpitations，因 cardiogenic shock 入院
    \item ECG：Atrial tachycardia 200 bpm
    \item Echo：Biventricular dysfunction，LVEDVi 117 mL/m$^2$，EF 12\%
    \item 治療：Impella CP $\rightarrow$ VA-ECMO (ECMELLA)
    \item RV EMB：無 myocardial inflammation，但有 early fibrosis，符合 dilated cardiomyopathy
    \item 成功 atrial arrhythmia ablation 後臨床改善
    \item 基因：Pathogenic TTN mutation
    \item 1 年後：LVEDVi 66 mL/m$^2$，EF 55\%
\end{itemize}

\subsection{Case 2：Myocarditis with Cardiogenic Shock}

\begin{itemize}
    \item 40 歲男性，retrosternal pain 及 flu-like symptoms 後出現 cardiogenic shock
    \item ECG：Sinus tachycardia with LBBB
    \item Echo：Moderate LV dilation，EF 25\%，moderate MR
    \item 治療：IABP + temporary pacemaker（因 AV block）
    \item LV EMB：Severe inflammation dominated by NK CD16+ lymphocytes；viral PCR negative
    \item 治療：Cortisone + azathioprine
    \item 15 天後出院，EF 正常
    \item CMR：Mild LV dilation，EF 53\%，midseptal intramural anterior LGE
\end{itemize}

\section{未來展望}

\begin{itemize}
    \item \textbf{Hybrid diagnostic protocols}：EMB 與 advanced imaging（CMR、scintigraphy、PET）整合
    \item \textbf{Real-time imaging guidance}：與 EAM systems 結合進行 targeted biopsies
    \item \textbf{Molecular and genetic analyses}：結合 tissue analysis 與 genetic testing 提供 personalized therapy
    \item \textbf{Gene therapy evaluation}：評估 cardiomyocyte transduction、histological features、protein expression
    \item \textbf{Training}：納入 interventional fellowship，simulation platforms 可增強訓練
    \item \textbf{Hub-and-spoke network}：建立專門中心網絡，優化 patient care
\end{itemize}

\section{結論}

EMB 對介入心臟科醫師而言是連結 invasive diagnostics 與 personalized care 的獨特機會。透過技術進步、與 non-invasive modalities 整合，可更精準有效地管理 myocarditis、cardiac amyloidosis、cardiac sarcoidosis 及 cardiomyopathies 等複雜心臟疾病。

\section{參考文獻}

\begin{enumerate}
    \item Fabris E, Porcari A, Bussani R, et al. Endomyocardial biopsy. \textit{EuroIntervention}. 2026;21:e19-e31.
    \item McDonagh TA, Metra M, Adamo M, et al. 2021 ESC Guidelines for the diagnosis and treatment of acute and chronic heart failure. \href{https://doi.org/10.1093/eurheartj/ehab368}{\textit{Eur Heart J}. 2021;42:3599-726.}
    \item Arbelo E, Protonotarios A, Gimeno JR, et al. 2023 ESC Guidelines for the management of cardiomyopathies. \href{https://doi.org/10.1093/eurheartj/ehad194}{\textit{Eur Heart J}. 2023;44:3503-626.}
    \item Drazner MH, Bozkurt B, Cooper LT, et al. 2024 ACC Expert Consensus Decision Pathway on Strategies and Criteria for the Diagnosis and Management of Myocarditis. \href{https://doi.org/10.1016/j.jacc.2024.10.080}{\textit{J Am Coll Cardiol}. 2025;85:391-431.}
\end{enumerate}

\end{document}
